\subsection{Universality}

\begin{definition}
For every $n \in \N$, Let $\SU(n)$ denote the special unitary group under matrix multiplication consisting of $n \times n$ unitary matrices with determinant 1. It is a normal subgroup of the unitary group $U(n)$, consisting of $n \times n$ unitary matrices.
\end{definition}

\begin{definition}
\label{def:universality}
A set of gates $\cS$ is (efficiently) \term{universal} if any $U \in \SU(2^n)$, for any $n \in \N$, that can be generated (in $\Poly(n)$ gates) by an arbitrary gate set $\cS'$, can also be generated by the gate set $\cS$ (in $\Poly(n)$ gates) \cite{MBQC-universal}. We say a set of gates $\cS$ is $d$-qubit universal if the above holds for $U \in \SU(2^d)$.
\end{definition}

In this thesis we will only focus on universality, not efficient universality. Thus, it may take exponential or more gates to simulate certain unitaries $U$. Also, in practice, the application of $U$ up to some measure and error $\epsilon$ suffices. This is called \term{approximate universality}. This thesis will focus solely on \term{exact universality}, where we desire no error, rather than approximate universality. Henceforth, all references to universality will refer to exact universality.

Furthermore, note that definition \ref{def:universality} is only concerned with implementing arbitrary unitaries, not necessarily arbitrary computation. That is, it does not specify whether the inputs and outputs are classical or quantum in nature. This matters because $\ket\phi = \ket0$ and $\ket\psi = -\ket0$ have the same measurement distribution and classical output, but $\ket\phi \neq \ket\psi$. This difference is called \term{global phase}, i.e. a scalar $e^{i\theta}$ multiplied to an entire state.

\begin{definition}
\ref{def:types-of-universal}
We use the following types of \term{universal computation} from \cite{MBQC-universal}:
\begin{itemize}
    \item CC-universal: For every $U \in \SU(n)$, the model can reproduce the distribution of local measurements performed on $U \{\ket{0^n}\}$.
    \item QC-universal: Given an arbitrary $n$-qubit state $\ket\phi$, for every $U \in \SU(2^n)$, the model can reproduce the distribution of local measurements performed on $U \ket\phi$.
    \item CQ-universal: For every $U \in SU(2^n)$, the model can produce the quantum state $\ket\phi = U\ket{0^n}$.
    \item QQ-universal: Given an arbitrary $n$-qubit state $\ket\phi$, for every $U \in \SU(2^n)$, the model can produce the quantum state $U \ket\phi$.
\end{itemize}
\end{definition}

These distinctions allow us to more finely categorize universality. For example, we may restate Lemma \ref{lemma:zx-ignore-scalar} as follows: Ignoring scalar factors during ZX rewrites, provided that the desired scalar is restored at the end, makes ZX-Calculus at most QC-universal, while keeping track of scalars throughout gives QQ-universality.

For this thesis, we will focus on (exact) CQ-universality, out of convenience and convention.\footnote{Many papers which examine the universality of MBQC resource states really only examine CQ-universality. For example, \cite{MBQC-universal}.} This distinction turns out to be quite important for algorithmic and implementation reasons, as we will see in Section \ref{sec:algorithmic-approaches}. Henceforth, universality shall refer to exact CQ-universality, unless otherwise noted.

\begin{fact}
\label{fact:1-qubit-and-entangling-universal}
Any 1-qubit universal gate set and a 2 qubit entangling gate is universal \cite{1qubit-entangling-universal}.
\end{fact}

\begin{proof}
The proof in \cite{1qubit-entangling-universal} turns out to be quite technical, involving many tricks from Lie Algebra. It surmounts to proving the following restatement of the theorem: Let $V$ be a 2-qubit gate. The following are equivalent:
\begin{enumerate}
    \item The collection of all 1-qubit gates $\cA$ together with $V$ is QQ-universal.
    \item The collection of all 1-qubit gates $\cA$ together with $V$ is exactly QQ-universal.
    \item $V$ is not primitive, or \term{imprimitive}.
\end{enumerate}
Here, say $V$ is primitive iff $V = S \otimes T$ or $V = (S \otimes T)P$ for some $1$-qubit gates $S$ and $T$, with $P$ being the swap gate.
\end{proof}

\begin{fact}
Any universal gate set must include at least one multi-qubit entangling gate, as it is impossible to entangle using only single-qubit gates.
\end{fact}

\begin{proof}
This follows from Fact \ref{fact:1-qubit-and-entangling-universal}.
\end{proof}